\documentclass[11pt,a4paper]{article}

\usepackage{amsmath,amssymb}
\usepackage[utf8]{inputenc}
\usepackage{hyperref}
\usepackage{geometry}
\usepackage{booktabs}
\usepackage{xcolor}
\usepackage{setspace}
\usepackage{enumitem}

\geometry{margin=1.15in}
\onehalfspacing

\newcommand{\RH}{\ensuremath{\mathrm{RH}}}

\title{%
  What Does a Machine Know When It Checks a Proof?\\[0.3cm]
  {\large Philosophical Implications of the Cathedral Formalization}%
}
\author{
  Jason Robert Gochanour
}
\date{April 19, 2026}

\begin{document}

\maketitle

\begin{abstract}
The \emph{Cathedral} is a machine-verified formalization in the Lean~4
theorem prover establishing that the Riemann Hypothesis is logically
equivalent to the decay of the Nyman--Beurling distance. It does not
prove the Riemann Hypothesis. It proves, with certainty beyond human
capacity for error, exactly what one \emph{would} need to prove.
This paper examines the philosophical questions raised by this
artifact: What is the epistemic status of a conditional proof? What
does it mean for a machine to ``catch'' a mathematical error that a
human would not? Can axioms serve as units of mathematical knowledge,
decomposing an open problem into precisely measured fragments of
remaining ignorance? We argue that the Cathedral represents a new
category of mathematical object---the \emph{certified reduction}---that
is neither a proof nor a conjecture, but a machine-witnessed
guarantee about the logical structure of an unsolved problem.
\end{abstract}

\tableofcontents

% ================================================================
\section{Introduction: A New Kind of Mathematical Object}
% ================================================================

On April 19, 2026, a 78-file computer program was compiled without
error. The program, written in the Lean~4 proof language and verified
against the Mathlib mathematical library, established the following:

\begin{quote}
\emph{The Riemann Hypothesis is true if and only if $d_N^2 \to 0$,}
\end{quote}

\noindent where $d_N^2$ is a concrete quantity measuring how well
certain sawtooth functions approximate a constant. The biconditional
is complete: both directions are proved. The proof depends on exactly
2 mathematical axioms (statements accepted without proof) plus 3
logical axioms built into Lean's type theory.

The Riemann Hypothesis remains unproved. But something has changed.
Before the Cathedral, the logical structure of the problem was
distributed across hundreds of papers, written over 167 years, in
inconsistent notation, with varying standards of rigor. Now it exists
as a single, machine-auditable artifact in which every deductive
step has been verified, every assumption is named, and the frontier
of human ignorance is precisely demarcated.

This artifact raises questions that mathematics alone cannot answer.
They are questions for philosophy.

% ================================================================
\section{The Epistemic Status of Conditional Proof}
\label{sec:conditional}
% ================================================================

\subsection{What Has Been Proved?}

The Cathedral proves a biconditional: $\RH \iff d_N^2 \to 0$. It
does \emph{not} prove either side independently. Neither the truth
nor the falsity of $\RH$ follows from the Cathedral alone.

This is not a new situation in mathematics. Many theorems are
conditional: ``If the Generalized Riemann Hypothesis is true,
then\ldots'' What \emph{is} new is the machine verification. The
biconditional is not a claim in a paper that might contain an error.
It is a statement that has been checked by a computer against a
formal system of logic (the Calculus of Inductive Constructions),
in which every step follows from the axioms by mechanical rule
application.

\begin{quote}
\emph{The conditional is certain. The condition is not.}
\end{quote}

This creates an unusual epistemic asymmetry. We know, with a
certainty that exceeds any human proof, that \RH{} is equivalent to
$d_N^2 \to 0$. But we do not know whether either statement is true.
The certainty attaches to the \emph{structure} of the implication,
not to its content.

\subsection{Is a Biconditional a ``Proof''?}

In the formalist tradition (Hilbert, Bourbaki), a proof is a finite
sequence of deductions from axioms. By this standard, the Cathedral
\emph{is} a proof---specifically, it is a proof of the proposition
``$\RH \iff d_N^2 \to 0$.'' The proposition happens to be a
biconditional rather than a simple statement, but formalism makes no
distinction between the two.

In the intuitionist tradition (Brouwer, Heyting), a proof of $P \iff Q$
is a pair of constructions: one transforming a proof of $P$ into a
proof of $Q$, and vice versa. The Cathedral provides exactly this:
the Lean file \texttt{Assembly/MainChain.lean} contains a term of
type $\RH \to (d_N^2 \to 0)$ and a term of type
$(d_N^2 \to 0) \to \RH$. (Lean~4 uses classical logic via
\texttt{Classical.choice}, so the constructions are not fully
constructive, but they are syntactically present.)

In the Lakatosian tradition (\emph{Proofs and Refutations}), a proof
is a social process of conjecture, criticism, and revision. By this
standard, the Cathedral is profoundly proof-like: it was built through
23 days of iterative refinement, with 96 archived files documenting
failed approaches, false lemmas, and superseded strategies. The
machine played the role of Lakatos's ``critic,'' catching errors that
the human would have missed (Section~\ref{sec:traps}).

\subsection{The Axiom as a Unit of Ignorance}

The most philosophically interesting feature of the Cathedral is its
treatment of axioms. In classical foundations (ZFC), axioms are the
starting points of all mathematics---they are accepted because they
are self-evident or because they codify our best intuitions about
sets.

In the Cathedral, axioms serve a fundamentally different purpose. They
are \emph{placeholders for unformalized knowledge}. Each axiom
represents a specific, named, precisely stated piece of mathematics
that is believed to be true (and has been proved by humans in
traditional papers) but has not yet been translated into Lean.

This makes the axiom a unit of \emph{measured ignorance}. The
Cathedral does not merely say ``we don't know how to prove \RH.''
It says: ``We don't know how to prove \RH, and here are the
\emph{exactly two} facts we would need to formalize in order to
reduce it to a question about $L^2$ approximation.''

\begin{quote}
\emph{An axiom is not a failure of proof. It is a unit of measurement
for what remains unproved.}
\end{quote}

The reduction from 56 axioms (March 27, Day 1) to 2 axioms (April 18,
Day 22) is not merely a technical achievement. It is an epistemic
one: the frontier of ignorance was compressed by 96\% in three weeks.

% ================================================================
\section{What Does the Machine ``Know''?}
\label{sec:knowledge}
% ================================================================

\subsection{Machine Verification Is Not Understanding}

The Lean type-checker that verifies the Cathedral does not
``understand'' the Riemann Hypothesis. It does not know what a prime
number is, what the zeta function does, or why the problem matters.
It checks that each step in the proof follows from the previous one
by the rules of the Calculus of Inductive Constructions. It performs
this check with perfect reliability, but without any semantic
comprehension.

This is a philosophical distinction with practical consequences.
A human mathematician who understands a proof can:

\begin{itemize}
  \item Generalize it to new contexts.
  \item Identify which steps are ``deep'' and which are routine.
  \item Predict where the proof technique might fail.
  \item Extract intuition for new conjectures.
\end{itemize}

The Lean type-checker can do none of these things. It can verify but
not understand, certify but not explain.

\subsection{But Verification Has Its Own Epistemic Value}

The flip side is that the machine verification provides something
that human understanding cannot: \emph{certainty}. Human proofs
contain errors. The history of mathematics is littered with proofs
that were accepted for years or decades before flaws were discovered
(Kempe's proof of the Four-Color Theorem, Wiles's initial proof of
Fermat's Last Theorem, Voevodsky's motivating experience with
homotopy groups).

The Cathedral's verification eliminates this class of error entirely.
If the Lean kernel is correct (a much smaller and more auditable
artifact than any mathematical proof), then the biconditional holds.
Period.

This creates a division of epistemic labor:

\begin{center}
\begin{tabular}{@{}ll@{}}
\toprule
\textbf{Machine provides} & \textbf{Human provides} \\
\midrule
Deductive certainty & Semantic understanding \\
Error elimination & Intuition and creativity \\
Axiom tracking & Mathematical strategy \\
Exhaustive checking & Selective attention \\
\bottomrule
\end{tabular}
\end{center}

Neither alone is sufficient. Together, they produce an artifact with
properties that neither could achieve independently.

\subsection{The AI Collaborator as Philosophical Agent}

The Cathedral was built not by a human alone, nor by a machine alone,
but by a human-AI pair. The AI assistant (an LLM-based coding agent)
wrote most of the Lean code, discovered Mathlib API functions,
and proposed proof strategies. The human directed the mathematical
architecture and made all strategic decisions.

This raises a question: \emph{who} proved the biconditional?

\begin{itemize}
  \item The human, who directed the strategy? But the human could not
    have written the 20,000 lines of Lean code alone.
  \item The AI, who wrote the code? But the AI did not understand
    the mathematics and could not have chosen the proof architecture.
  \item The Lean kernel, which verified the proof? But the kernel
    made no creative contribution---it only checked.
\end{itemize}

Perhaps the correct answer is that mathematical proof is becoming
a \emph{collective} activity in a new sense: not a team of humans,
but a team of heterogeneous cognitive agents, each contributing a
different epistemic capacity. The result is a proof that no single
agent could have produced.

% ================================================================
\section{Five Traps and the Limits of Informal Reasoning}
\label{sec:traps}
% ================================================================

During construction, the machine-verification process detected five
mathematical errors that would have been invisible to informal proof.
Each illuminates a different limitation of human mathematical
reasoning:

\subsection{The Basis Trap}

Two natural choices for the sawtooth basis---$\{k/x\}$ and
$\{1/(kx)\}$---appear interchangeable to informal reasoning.
They are not. The first generates $L^2(0,1)$ unconditionally
(independent of \RH), making any approximation result trivially
true. The second has approximation quality governed by the zeta zeros.

An informal proof using the wrong basis would appear correct at
every step. No referee would catch it, because the error is in the
\emph{choice of definition}, not in any deductive step. The machine
caught it because the Lean type system enforces that the basis
functions must match the type signature of the Nyman--Beurling
criterion.

\textbf{Philosophical lesson:} Definitions are not epistemically
neutral. The choice of how to formalize a concept constrains which
theorems can be proved about it. Machines enforce definitional
consistency in a way that human attention cannot.

\subsection{The Cancellation Trap}

The triangle inequality gives $\|1 - f\|_2 \leq 1 + \|f\|_2$,
yielding $d_N^2 \leq 4$ for a quantity approaching 0. The bound is
mathematically correct but useless, because it destroys the
cancellation $1 - 2b^Tv + v^TGv \to 0$ that makes $d_N^2$ small.

An informal proof might invoke the triangle inequality without
noticing that it loses the essential structure. The machine caught
this because the Lean proof could not close the gap between the
bound and the target.

\textbf{Philosophical lesson:} Not all valid inequalities are useful.
Some destroy the very structure they are meant to control. This is
invisible to syntactic reasoning but immediately apparent to a
type-checker that demands a complete chain from hypothesis to
conclusion.

\subsection{The Ghost Axiom Trap}

Nine mathematical propositions existed simultaneously as axioms
(unproved) and theorems (proved) in different files. The module
system allowed both to coexist without contradiction. Downstream
theorems could silently depend on the axiom rather than the theorem,
inflating the proof's apparent weakness.

\textbf{Philosophical lesson:} In modular epistemic systems, the same
proposition can have different justification statuses in different
contexts. Coherence checking across modules is a non-trivial epistemic
task that humans perform poorly and machines perform well.

% ================================================================
\section{The Cathedral as Philosophical Object}
\label{sec:object}
% ================================================================

\subsection{A Certified Reduction}

We propose that the Cathedral represents a new category of
mathematical object: the \emph{certified reduction}. A certified
reduction has three components:

\begin{enumerate}
  \item A \textbf{target proposition} (here: \RH).
  \item A \textbf{reduced proposition} (here: $d_N^2 \to 0$).
  \item A \textbf{machine-verified biconditional} connecting them,
    with every assumption explicitly named.
\end{enumerate}

A certified reduction is stronger than a conjecture (it establishes
logical equivalence) but weaker than a proof (it does not resolve
either proposition). Its epistemic value lies in the precision of
the residual: the named axioms constitute an \emph{exact measure}
of what remains unknown.

\subsection{Axiom Reduction as Epistemic Progress}

In traditional mathematics, progress on an open problem is measured
by theorems: new results that narrow the gap between what is known
and what is conjectured. The Cathedral suggests a complementary
measure: \emph{axiom count on the critical path}.

\begin{center}
\begin{tabular}{@{}cl@{}}
\toprule
\textbf{Axioms} & \textbf{Interpretation} \\
\midrule
56 & ``We have a rough sketch and many assumptions.'' \\
12 & ``We have a solid architecture with specific gaps.'' \\
5 & ``We can name every remaining piece.'' \\
2 & ``Two facts stand between us and certified equivalence.'' \\
0 & ``The equivalence is fully machine-verified.'' \\
\bottomrule
\end{tabular}
\end{center}

The progression from 56 to 2 is epistemic progress, even though
neither \RH{} nor its negation was established. The frontier of
ignorance shrank. The remaining gaps were named, located, and
measured. This is a form of understanding that traditional
mathematics---which classifies results as ``proved'' or ``open''
with no gradations---does not capture.

\subsection{The Archaeology of Failed Proofs}

The Cathedral's archive of 96 superseded files is, to our knowledge,
unprecedented in mathematical practice. Mathematicians typically
discard failed approaches. The Cathedral preserves them, creating a
\emph{fossil record} of mathematical exploration.

This archive has philosophical significance. It makes visible the
\emph{negative knowledge} that is normally invisible: approaches that
don't work, lemmas that are false, definitions that lead nowhere.
Lakatos argued that refutations are as important as proofs. The
Cathedral's archive is a concrete embodiment of this principle.

The 54\% discard rate (96 archived files out of 174 total) quantifies
something that mathematics has never measured before: the ratio of
exploration to exploitation in a proof search. More than half of all
mathematical effort was invested in paths that were eventually
abandoned. This is not failure; it is the cost of discovery.

% ================================================================
\section{Implications for Mathematical Practice}
\label{sec:practice}
% ================================================================

\subsection{The End of the Informal Era?}

The Cathedral was not built because informal mathematics is inadequate
for the Nyman--Beurling criterion. Every theorem in the Cathedral
could, in principle, be written as a traditional paper proof. The
machine verification adds no new mathematical content.

What it adds is \emph{trust infrastructure}. In a world where
LLMs can generate plausible-sounding but incorrect mathematical
proofs, machine verification provides a trust anchor that is
independent of human judgment, reputation, or social consensus.

\begin{quote}
\emph{Lean does not care who you are. It cares whether your proof
type-checks.}
\end{quote}

This is both a democratization (anyone can submit a machine-verified
proof) and a raising of standards (no amount of prestige can override
a type error).

\subsection{The Decomposition of Mathematical Problems}

If axioms can serve as units of ignorance, then open problems can
be \emph{decomposed} into precisely stated sub-problems, each
corresponding to an axiom. This transforms the sociology of
mathematical research: instead of one person (or group) attempting
an entire proof, the problem can be distributed.

The Cathedral's 2 remaining axioms are, in effect, \emph{bounty
postings}: precisely stated mathematical challenges, with a
machine-verified guarantee that formalizing them would complete the
reduction. This is a new mode of mathematical collaboration, enabled
by the precision of formal systems.

\subsection{What Counts as a Contribution?}

Traditional mathematics values theorems. A mathematician's
contribution is measured by the propositions they prove. The
Cathedral suggests a broader accounting:

\begin{itemize}
  \item \textbf{Axiom elimination} (proving a previously assumed result)
    is a contribution.
  \item \textbf{Architecture design} (choosing the right decomposition
    of a problem into modules) is a contribution.
  \item \textbf{Bug detection} (finding the basis trap, the cancellation
    trap) is a contribution.
  \item \textbf{Archival} (documenting failed approaches so others
    don't repeat them) is a contribution.
\end{itemize}

None of these fit neatly into the ``theorem proved'' metric. But all
of them advanced mathematical knowledge in a measurable way.

% ================================================================
\section{The Question of Beauty}
\label{sec:beauty}
% ================================================================

Mathematicians often speak of beauty in proofs. Hardy's \emph{A
Mathematician's Apology} elevates aesthetic criteria to central
importance. The Cathedral raises the question: can a machine-verified
proof be beautiful?

The Rank-1 Mellin Miracle is, by any standard, an elegant result.
At a zeta zero $\rho$, the Mellin transform of the B\'aez-Duarte
basis function collapses to $\mathcal{M}[h_k](\rho) = 1/(k(\rho-1))$.
The entire infinite-dimensional computation reduces to a single
multiplication. The $k$-dependence and $\rho$-dependence factorize
completely.

This is beautiful. And it was discovered during machine verification,
when a general strategy (Hahn--Banach separation) was replaced by a
specific computation that turned out to have unexpected structure.
The machine did not appreciate the beauty. The human did. But the
beauty would not have been found without the machine's insistence
on explicit computation.

\begin{quote}
\emph{The machine demands rigor. The human supplies meaning.
Neither alone produces beauty. Together, sometimes, they do.}
\end{quote}

% ================================================================
\section{Conclusion: The Cathedral as Epistemological Experiment}
% ================================================================

The Cathedral is many things: a 78-file Lean program, a reduction
of the Riemann Hypothesis, a case study in proof engineering, a
record of human-AI collaboration. But above all, it is an
\emph{epistemological experiment}: an attempt to make the structure
of mathematical knowledge visible, precise, and machine-auditable.

The experiment yields several findings:

\begin{enumerate}
  \item \textbf{Conditional proofs have precise epistemic value},
    measurable by the number and nature of their axioms.
  \item \textbf{Machine verification and human understanding are
    complementary}, not competitive, epistemic capacities.
  \item \textbf{Axioms can serve as units of measured ignorance},
    decomposing open problems into named, located, precisely
    stated fragments.
  \item \textbf{Failed approaches have epistemic value} when
    archived systematically.
  \item \textbf{Mathematical beauty can emerge from machine-assisted
    exploration}, though it requires human recognition to be
    appreciated.
\end{enumerate}

Two axioms remain. The Riemann Hypothesis is neither proved nor
disproved. But the space of possible proofs has been mapped with
unprecedented precision, and that mapping is itself a form of
mathematical knowledge---one that our philosophical vocabulary is
only beginning to accommodate.

\vspace{1cm}
\begin{flushright}
\emph{The Cathedral does not answer the question.}\\
\emph{It makes the question precise.}\\
\emph{That is philosophy's oldest aspiration.}
\end{flushright}

\vspace{1cm}

\begin{thebibliography}{99}

\bibitem{lakatos}
I.~Lakatos, \emph{Proofs and Refutations: The Logic of Mathematical
Discovery}, Cambridge University Press, 1976.

\bibitem{hardy}
G.H.~Hardy, \emph{A Mathematician's Apology}, Cambridge University
Press, 1940.

\bibitem{voevodsky}
V.~Voevodsky, \emph{The Origins and Motivations of Univalent
Foundations}, IAS, 2014.

\bibitem{avigad}
J.~Avigad, \emph{The mechanization of mathematics}, Notices of the
AMS \textbf{65} (2018), 681--690.

\bibitem{appel1977}
K.~Appel and W.~Haken, \emph{Every planar map is four colorable},
Illinois J. Math. \textbf{21} (1977), 429--490.

\bibitem{buzzard}
K.~Buzzard, \emph{The rise of formalism in mathematics},
ICM 2022 proceedings.

\end{thebibliography}

\end{document}
